% Recovery of 02_trit_geometrias and holonomy with memory.
\subsubsection{Temporal parametrization of compatibility and tritic orientation}
\label{trit-temporal-compatibilidad}

The temporal interpretation of TRIT is established on transitions of the
complete state. Let \(x\xrightarrow{\rm adm}y\) be the admissibility relation
determined by TPK operations on \(\widehat X\). A trajectory parametrized
by a discretely ordered set \(\mathcal T\) is a map
\[
 \gamma:\mathcal T\longrightarrow\widehat X,\qquad
 \gamma(t)\xrightarrow{\rm adm}\gamma(t^+),
\]
where \(t^+\) is the successor of \(t\). The compatibility relation
precedes the choice of parameter: the order allows its transitions to
be expressed as an evolution. When the term section is used, the
projection \(p:E\to\mathcal T\) and condition
\(p\circ\gamma=\operatorname{id}_{\mathcal T}\) are also declared.
The state space, temporal order, and transported information are thereby specified.

The signature \(\tau\) selects a quadratic regime and admits a temporal
reading relative to this parametrization. Type \(+1\), with
\(J_{+1}^{2}=-I\), represents return and memory; type \(0\), with
\(J_0^{2}=0\), represents the transition threshold; type \(-1\), with
\(J_{-1}^{2}=I\), represents opening in two eigendirections. In the
model's temporal terminology, these functions correspond, respectively,
to past, present, and future. The correspondence concerns the operational
position of a transition within a history: retained antecedent, update,
and admissible prolongation.

This interpretation acquires content through three distinct operations.
Restriction of a trajectory recovers the prefix already constructed.
The update \(\mathcal U_t\) composes selection, transport, and recording
in the observed transition. Prolongation considers continuations that
preserve the boundary conditions of that prefix. If
\(\mathcal P_n(x)\) denotes the set of admissible continuations of
length \(n\) from \(x\), deleting the last step defines
\[
 r_n:\mathcal P_{n+1}(x)\longrightarrow\mathcal P_n(x).
\]
An unlimited continuation is a collection \((\gamma_n)_{n\geq1}\)
satisfying \(r_n(\gamma_{n+1})=\gamma_n\). The inverse system thus
preserves the restrictions each depth imposes on preceding depths.
Existence of a compatible collection is established in the corresponding
prolongation domain; merely specifying the sets \(\mathcal P_n(x)\)
does not replace that proof.

Bidirectionality therefore combines reconstruction of antecedents with
compatibility of continuations. On the signature sequence, the involution
\(C_{\rm rev}\) reverses order and changes the sign of each component.
The trajectory is also subject to the transformations of position, sheet,
and boundary specified by transport. The reversed signature is an
observation of the conjugate trajectory, not a replacement for its other
coordinates. The threshold remains distinguished: the visible value
\(0\) retains its transition function even when its quotient or
historical record is nonzero.

The relation with exponential geometries is equally precise. In each
fixed regime, \(\exp(tJ_\tau)\) composes parameters and preserves the
algebraic norm. Elliptic closure belongs to that local representation.
Phase return of \(\Gamma_9\), by contrast, acts on the trajectory
with memory, according to \eqref{eq:holonomia-memoria}. A trajectory
can recover its phase while retaining a new antecedent: historical data
have changed although the phase observation is the same. The parabolic
threshold retains a linear displacement; hyperbolic opening separates
two eigendirections of expansion and contraction. These three
realizations provide composition laws, whereas the TPK determines their
selection and effective concatenation.

The aperiodic suspension in Section \ref{sec:generacion} will make this
distinction explicit through a finite phase factor, an irrational rotation,
and a memory degree. The link between TRIT and temporality is thus
expressed by operations on trajectories and their quadratic representations.
Its content exceeds a classification of signs: it preserves the regime,
the direction of transport, and the relation between antecedent,
transition, and continuation.
